% Capítulo — TCC 1 (substituído pela Conclusão no TCC 2)
\chapter{Considerações finais do TCC 1}
\label{chap:consideracoes}

Esta primeira etapa definiu o problema, o princípio de projeto e o desenho da
avaliação. O problema é a descoberta de chaves de integração entre bases
governamentais que descrevem as mesmas entidades com nomes, formatos e
códigos diferentes: no Gov Hub, 28 das 29 condições de junção entre fontes
precisam de alguma transformação. O princípio é a agência limitada: o LLM
contribui com conhecimento de domínio, mas escolhe dentro de um catálogo
fechado ou se abstém, e código determinístico constrói e verifica a junção. O
mesmo princípio limita o agente de IA que constrói o sistema, por meio de
especificações revisadas e testes escondidos dele.

Se essa limitação melhora a acurácia em relação ao \emph{matching} clássico e
ao uso livre do LLM é uma questão empírica, assim como saber se uma
especificação completa melhora o trabalho do agente. O estudo foi desenhado
para responder às duas. No MinC, o sistema será aplicado a fontes que o
ministério já ingere mas ainda não cruza, com a meta de entregar novos
cruzamentos validados.

Para o TCC~2, os próximos passos são, em ordem:

\begin{enumerate}
\item construir por SDD a chamada de decisão tipada, a abstenção, o catálogo,
o achatamento de JSON e a renderização dbt;
\item montar o benchmark (pares reais com \emph{snapshots}, gerador de pares
perturbados e divisão sem vazamento), o \emph{harness} e os
\emph{baselines}, e escrever o Mapa de Fluxo;
\item congelar o conjunto de teste, o catálogo e o Mapa de Fluxo;
\item executar a QP1 a QP3 e aplicar o sistema aos pares do MinC (QP4);
\item consolidar os registros da QP5 e escrever os capítulos de resultados e
discussão.
\end{enumerate}

O sistema, o benchmark (gabaritos e gerador de pares) e os artefatos do SDD
serão publicados no repositório do projeto.
